\documentclass{scrartcl}
\usepackage{graphicx}  % required for inserting images

% for cyrillic and Bulgarian language
\usepackage[utf8]{inputenc}
\usepackage[T2A]{fontenc}
\usepackage[bulgarian]{babel}

% math packages
\usepackage{amsmath}
\usepackage{float}
\usepackage{amsfonts}
\usepackage{amssymb}
\usepackage{amsthm}
\usepackage{cancel}

\usepackage{ragged2e}  % adds FlushLeft

\usepackage{listings}
\usepackage{xcolor}

\definecolor{codegreen}{rgb}{0,0.6,0}
\definecolor{codegray}{rgb}{0.5,0.5,0.5}
\definecolor{codepurple}{rgb}{0.58,0,0.82}
\definecolor{backcolour}{rgb}{0.95,0.95,0.92}

\lstdefinestyle{mystyle}{
    language=Octave,
    backgroundcolor=\color{backcolour},   
    commentstyle=\color{codegreen},
    keywordstyle=\color{magenta},
    numberstyle=\tiny\color{codegray},
    stringstyle=\color{codepurple},
    basicstyle=\ttfamily\footnotesize,
    breakatwhitespace=false,         
    breaklines=true,
    captionpos=b,
    keepspaces=false,                 
    numbersep=5pt,
    showspaces=false,
    showstringspaces=false,
    showtabs=false,
    tabsize=2,
    columns=fullflexible,
}

\lstset{style=mystyle}

\usepackage[colorlinks=true, linkcolor=blue, urlcolor=cyan]{hyperref}
\usepackage{tikz}
\usetikzlibrary{arrows.meta}

\newtheorem{theorem}{Теорема}
\newtheorem{definition}{Дефиниция}

\newtheorem{lemma}{Лема}
\newtheorem{corollary}{Следствие}
\newtheorem{remark}{Забележка}

\title{Вероятности. Основни разпределения.}
\author{Ивайло Андреев + gemini 3 Flash, Claude Sonnet 4.6}
\date{14 юни 2026}

\begin{document}
\maketitle

\tableofcontents

\newpage

\begin{FlushLeft}

\section*{Полезни формули}

\begin{enumerate}

\item \textbf{Формула за бинома на Нютон}
\[
(a+b)^n
=
\sum_{k=0}^{n}
\binom{n}{k}
a^k b^{n-k}
\]

\item \textbf{Сума на първите $n$ естествени числа}
\[
\sum_{k=1}^{n} k
=
\frac{n(n+1)}{2}
\]

\item \textbf{Сума на квадратите на първите $n$ естествени числа}
\[
\sum_{k=1}^{n} k^2
=
\frac{n(n+1)(2n+1)}{6}
\]

\item \textbf{Безкрайна геометрична прогресия}
\[
\sum_{k=0}^{\infty} x^k
=
\frac{1}{1-x},
\qquad |x|<1
\]

\item \textbf{Ред на Маклорен за експоненциалната функция}
\[
e^x
=
\sum_{k=0}^{\infty}
\frac{x^k}{k!}
\]

\end{enumerate}

\section{Основни понятия и свойства на вероятностите}

\subsection{Вероятностно пространство и свойства на вероятността}
Нека $\Omega$ е непразно множество от всички възможни елементарни изходи от даден случаен експеримент, а $\mathcal{A}$ е $\sigma$-алгебра от подмножества на $\Omega$ (наречени събития). 

\begin{definition}
Функцията $\mathbb{P}: \mathcal{A} \to [0, 1]$ се нарича \textbf{вероятност}, ако удовлетворява аксиомите на Колмогоров:
\begin{enumerate}
    \item \textbf{Неотрицателност}: $\mathbb{P}(A) \ge 0$ за всяко събитие $A \in \mathcal{A}$.
    \item \textbf{Нормираност}: Вероятността на сигурното събитие $\Omega$ е равна на единица: $\mathbb{P}(\Omega) = 1$.
    \item \textbf{Счетна адитивност}: За всяка поредица от взаимно изключващи се събития $A_1, A_2, \dots$ ($A_i \cap A_j = \emptyset$ при $i \neq j$) е изпълнено:
    \[ \mathbb{P}\left(\bigcup_{i=1}^{\infty} A_i\right) = \sum_{i=1}^{\infty} \mathbb{P}(A_i) \]
\end{enumerate}
\end{definition}

\begin{definition}
    Вероятностно пространство наричаме наредената тройка $(\Omega, \mathcal{A}, \mathbb{P})$.
\end{definition}

Вероятностната мярка има следните основни свойства:

\begin{itemize}
    \item \textbf{Адитивност (крайна)}: Ако $A \cap B = \emptyset$, то $\mathbb{P}(A \cup B) = \mathbb{P}(A) + \mathbb{P}(B)$.
    \item \textbf{Монотонност}: Ако събитието $A$ е подмножество на събитието $B$ ($A \subseteq B$), то $\mathbb{P}(A) \le \mathbb{P}(B)$. Това следва от представянето $B = A \cup (B \setminus A)$, където $A \cap (B \setminus A) = \emptyset$, откъдето $\mathbb{P}(B) = \mathbb{P}(A) + \mathbb{P}(B \setminus A) \ge \mathbb{P}(A)$.
    \item \textbf{Вероятност на допълнително събитие}: $\mathbb{P}(A^c) = 1 - \mathbb{P}(A)$.
    \item \textbf{Вероятност на невъзможното събитие}: $\mathbb{P}(\emptyset) = 0$.
\end{itemize}

\subsection{Дефиниция на дискретно вероятностно разпределение}
Случайна величина $X$ е измерима функция, дефинирана над пространството от елементарни събития, $X: \Omega \to \mathbb{R}$, ако за всяко $a \in \mathbb{R}$ е изпълнено, че $X^{-1}((-\infty, a)) \in \mathcal{A}$.

\begin{definition}
Казваме, че случайната величина $X$ е \textbf{дискретна}, ако нейното множество от стойности $X(\Omega) = \{x_1, x_2, \dots, x_n, \dots\}$ е крайно или изброимо. \textbf{Дискретно вероятностно разпределение} на $X$ се нарича съвкупността от вероятностите $\{p_i\}$, с които величината приема своите стойности:
\[ p_i = \mathbb{P}(X = x_i), \quad i = 1, 2, \dots \]
\end{definition}

Дискретната \textbf{функция на разпределение} на $X$ се дефинира като:
\[ F_X(x) = \mathbb{P}(X \le x) = \sum_{x_i \le x} p_i \]

За да бъде едно разпределение коректно дефинирано, вероятностите $\{p_i\}$ задължително трябва да изпълняват следните две свойства (свойства на разпределението):
\begin{enumerate}
    \item \textbf{Неотрицателност}: $p_i \ge 0$ за всяко $i$.
    \item \textbf{Условие за нормировка}: Сумата от вероятностите на всички възможни изходи трябва да е единица:
    \[ \sum_{i} p_i = 1 \]
\end{enumerate}

Често представяме дискретните случайни величини със следната таблица:

\begin{tabular}{|c|c|c|c|}
  \hline
  $X$ & $x_1$ & $x_2$ & $\dots$ \\ \hline
  $\mathbb{P}(X = x_i)$ & $p_1$ & $p_2$ & $\dots$ \\ \hline
\end{tabular}

\subsection{Моменти на случайни величини}
\begin{definition}
\textbf{Математическо очакване} (първи начален момент) на дискретна случайна величина $X$ се нарича числото:
\[ \mathbb{E}X = \sum_{i} x_i p_i \]
дефинирано при условие, че редът е абсолютно сходящ ($\sum_{i} |x_i| p_i < \infty$).
\end{definition}

Очакването е линеен оператор относно случайни величини. Това означава, че ако $X, Y$ са (дискретни) случайни величини и $\alpha, \beta, \gamma$ са реални скалари, то е изпълнено равенството:

$$\mathbb{E}(\alpha X + \beta Y + \gamma) = \alpha\mathbb{E}X + \beta\mathbb{E}Y + \gamma$$

\begin{definition}
\textbf{Дисперсия} (втори централен момент) на случайна величина $X$ измерва разсейването на стойностите около тяхното математическо очакване и се дефинира като:
\[ \mathbb{D}X = \mathbb{E}(X - \mathbb{E}X)^2 = \sum_{i} (x_i - \mathbb{E}X)^2 p_i \]
\end{definition}
За практически пресмятания често се използва по-удобната еквивалентна формула:
\[ \mathbb{D}X = \mathbb{E}X^2 - (\mathbb{E}X)^2, \quad \text{където } \mathbb{E}X^2 = \sum_{i} x_i^2 p_i \]

В общия случай дисперсията не е линеен оператор относно случайни величини.
Ако $X$ и $Y$ са независими случайни величини, т.е. за всякакви стойности $k_1$ и $k_2$ е изпълнено
\[
\mathbb{P}(X = k_1,\, Y = k_2) = \mathbb{P}(X = k_1)\,\mathbb{P}(Y = k_2),
\]
то тогава е изпълнено равенството:

$$\mathbb{D}(X + Y) = \mathbb{D}X + \mathbb{D}Y$$

\section{Пораждащи функции и техни свойства}

При аналитичното изследване на дискретни случайни величини, приемащи цели неотрицателни стойности ($X \in \{0, 1, 2, \dots\}$), мощен инструмент се явява вероятностната пораждаща функция.

\begin{definition}
\textbf{Пораждаща функция} на дискретна случайна величина $X$ се нарича функцията на комплексния или реалния аргумент $s$, дефинирана като математическото очакване на $s^X$:
\[ \varphi_X(s) = \mathbb{E}[s^X] = \sum_{k=0}^{\infty} p_k s^k, \quad \text{където } p_k = \mathbb{P}(X = k) \]
\end{definition}

\subsection{Основни свойства на пораждащите функции}
\begin{enumerate}
    \item \textbf{Сходимост}: Тъй като $\sum_{k=0}^{\infty} p_k = 1$, редът, дефиниращ $\varphi_X(s)$, е абсолютно сходящ в затворения кръг $|s| \le 1$.
    \item \textbf{Стойност в единицата}: $\varphi_X(1) = \sum_{k=0}^{\infty} p_k = 1$.
    \item \textbf{Връзка с моментите (намиране на математическо очакване)}: Първата производна на пораждащата функция в точката $s=1$ дава математическото очакване на случайната величина:
    \[ \varphi_X'(1) = \left. \sum_{k=1}^{\infty} k p_k s^{k-1} \right|_{s=1} = \sum_{k=1}^{\infty} k p_k = \mathbb{E}X \]
    \item \textbf{Връзка с моментите (намиране на дисперсия)}: Втората производна в точката $s=1$ дава втория факториален момент:
    \[ \varphi_X''(1) = \left. \sum_{k=2}^{\infty} k(k-1) p_k s^{k-2} \right|_{s=1} = \mathbb{E}[X(X-1)] = \mathbb{E}X^2 - \mathbb{E}X \]
    Следователно дисперсията може да се изрази чрез производните в точката 1:
    \[ \mathbb{D}X = \varphi_X''(1) + \varphi_X'(1) - (\varphi_X'(1))^2 \]
    \item \textbf{Пораждаща функция на сума от независими случайни величини}: Ако $X$ и $Y$ са независими случайни величини, то пораждащата функция на тяхната сума е произведение от техните пораждащи функции:
    \[ \varphi_{X+Y}(s) = \varphi_X(s) \cdot \varphi_Y(s) \]
\end{enumerate}

\section{Анализ на основните дискретни разпределения}

\subsection{Дискретно равномерно разпределение}
\subsubsection{Модел и дефиниция}
Нека $X$ е случайна величина, описваща резултата от случаен избор на един елемент измежду $n$ еднакво вероятни изхода. Тогава $X$ има дискретно равномерно разпределение $X \sim U\{1, 2, \dots, n\}$, зададено чрез:
\[ \mathbb{P}(X = k) = \frac{1}{n}, \quad k = 1, 2, \dots, n \]

\textbf{Пример}: Хвърляне на стандартен симетричен куб с шест страни. Случайната величина $X$, съпоставяща броя на падналите точки, има равномерно разпределение с $n = 6$.

\subsubsection{Пораждаща функция и моменти}
Намираме пораждащата функция:
\[ \varphi_X(s) = \mathbb{E}[s^X] = \frac{1}{n} \sum_{k=1}^{n} s^k = \frac{s(s^n - 1)}{n(s - 1)}, \quad s \neq 1 \]
Математическото очакване намираме директно от дефиницията, използвайки формулата за сума на първите $n$ естествени числа ($\sum_{k=1}^n k = \frac{n(n+1)}{2}$):
\[ \mathbb{E}X = \frac{1}{n} \sum_{k=1}^{n} k = \frac{1}{n} \cdot \frac{n(n+1)}{2} = \frac{n+1}{2} \]
Вторият момент намираме директно от дефиницията, използвайки формулата за сума на квадратите ($\sum_{k=1}^n k^2 = \frac{n(n+1)(2n+1)}{6}$):
\[ \mathbb{E}X^2 = \frac{1}{n} \sum_{k=1}^{n} k^2 = \frac{(n+1)(2n+1)}{6} \]
\[ \mathbb{D}X = \mathbb{E}X^2 - (\mathbb{E}X)^2 = \frac{(n+1)(2n+1)}{6} - \left(\frac{n+1}{2}\right)^2 = \frac{n^2 - 1}{12} \]

\textit{Забележка}: В литературата често се среща еквивалентен вариант на разпределението, получен чрез линейна промяна на случайната величина. Ако $X \sim U\{1, \dots, n\}$, то $X' = X + m - 1 \sim U\{m, \dots, m+n-1\}$.

\subsection{Биномно разпределение}
\subsubsection{Модел и дефиниция}
Нека $X$ е случайна величина, описваща броя на успехите в $n$ независими опита от схемата на Бернули с вероятност за успех $p$ и неуспех $q = 1 - p$. Тогава $X$ има биномно разпределение $X \sim Bi(n, p)$, зададено чрез:
\[ \mathbb{P}(X = k) = \binom{n}{k} p^k q^{n-k}, \quad k = 0, 1, \dots, n \]
Условието за нормировка е изпълнено по формулата на бинома на Нютон: $\sum_{k=0}^n \binom{n}{k} p^k q^{n-k} = (p+q)^n = 1$.

\textbf{Пример}: Тестване на партида от 100 детайла, произведени на автоматична линия, при която вероятността за дефект на всеки детайл е $0.05$. Случайната величина $X$ – брой дефектни детайли, е $X \sim Bi(100,\, 0.05)$.

\subsubsection{Пораждаща функция и моменти}
Намираме пораждащата функция:
\[ \varphi_X(s) = \sum_{k=0}^{n} \binom{n}{k} (ps)^k q^{n-k} = (ps + q)^n \]
Намираме първата и втората производна по $s$:
\[ \varphi_X'(s) = np(ps + q)^{n-1}, \qquad \varphi_X''(s) = n(n-1)p^2(ps + q)^{n-2} \]
Заместваме $s = 1$ (при $p + q = 1$):
\[ \mathbb{E}X = \varphi_X'(1) = np \]
\[ \mathbb{D}X = \varphi_X''(1) + \varphi_X'(1) - (\varphi_X'(1))^2 = n(n-1)p^2 + np - n^2p^2 = np(1-p) = npq \]

\subsection{Геометрично разпределение}
\subsubsection{Модел и дефиниция}
Нека $X$ е случайна величина, описваща броя на неуспехите преди настъпването на първия успех в независими опити от схемата на Бернули с вероятност за успех $p$ и неуспех $q = 1 - p$. Тогава $X$ има геометрично разпределение $X \sim Ge(p)$, зададено чрез:
\[ \mathbb{P}(X = k) = q^k p, \quad k = 0, 1, 2, \dots \]
Условието за нормировка е изпълнено по формулата за безкрайна геометрична прогресия: $\sum_{k=0}^{\infty} q^k p = \frac{p}{1 - q} = 1$.

\textbf{Пример}: Баскетболист стреля към коша с вероятност за попадение $p$ при всеки независим изстрел. Случайната величина $X$ – брой на пропуските преди първото попадение, е $X \sim Ge(p)$.

\subsubsection{Пораждаща функция и моменти}
Намираме пораждащата функция ($\sum_{k=0}^{\infty} x^k = \frac{1}{1-x}$ за $|x| < 1$):
\[ \varphi_X(s) = \sum_{k=0}^{\infty} q^k p \cdot s^k = p \sum_{k=0}^{\infty} (qs)^k = \frac{p}{1 - qs}, \quad |qs| < 1 \]
Намираме производните по $s$:
\[ \varphi_X'(s) = \frac{pq}{(1 - qs)^2}, \qquad \varphi_X''(s) = \frac{2pq^2}{(1 - qs)^3} \]
Заместваме $s = 1$:
\[ \mathbb{E}X = \varphi_X'(1) = \frac{pq}{p^2} = \frac{q}{p} \]
\[ \mathbb{D}X = \varphi_X''(1) + \varphi_X'(1) - (\varphi_X'(1))^2 = \frac{2q^2}{p^2} + \frac{q}{p} - \frac{q^2}{p^2} = \frac{q^2 + pq}{p^2} = \frac{q}{p^2} \]

\textit{Забележка}: В литературата често се среща еквивалентен вариант на разпределението, получен чрез линейна промяна на случайната величина. Ако $X \sim Ge(p)$ брои неуспехите преди първия успех, то $X' = X + 1$ брои общия брой опити до (включително) първия успех.

\subsection{Поасоново разпределение}
\subsubsection{Модел и дефиниция}
Нека $X$ е случайна величина, описваща броя на случайни независими събития, настъпващи с постоянна средна интензивност $\lambda > 0$ в дадена единица (период от време, дължина, площ и др.). Тогава $X$ има Поасоново разпределение $X \sim Po(\lambda)$, зададено чрез:
\[ \mathbb{P}(X = k) = \frac{\lambda^k e^{-\lambda}}{k!}, \quad k = 0, 1, 2, \dots \]
Условието за нормировка е изпълнено по разлагането на $e^\lambda$ в ред на Маклорен: $\sum_{k=0}^{\infty} \frac{\lambda^k e^{-\lambda}}{k!} = e^{-\lambda} e^{\lambda} = 1$.

\textbf{Пример}: Брой на телефонни повиквания, пристигащи към сървър в рамките на фиксиран интервал от една минута, ако те настъпват независимо и хомогенно. Разпределението е $Po(\lambda)$, където $\lambda$ е средният брой повиквания на минута.

\subsubsection{Пораждаща функция и моменти}
Намираме пораждащата функция:
\[ \varphi_X(s) = e^{-\lambda} \sum_{k=0}^{\infty} \frac{(\lambda s)^k}{k!} = e^{-\lambda} \cdot e^{\lambda s} = e^{\lambda(s - 1)} \]
Намираме производните по $s$:
\[ \varphi_X'(s) = \lambda e^{\lambda(s-1)}, \qquad \varphi_X''(s) = \lambda^2 e^{\lambda(s-1)} \]
Заместваме $s = 1$:
\[ \mathbb{E}X = \varphi_X'(1) = \lambda \]
\[ \mathbb{D}X = \varphi_X''(1) + \varphi_X'(1) - (\varphi_X'(1))^2 = \lambda^2 + \lambda - \lambda^2 = \lambda \]

\textit{Забележка}: Свойство на Поасоновото разпределение е, че $\mathbb{E}X = \mathbb{D}X = \lambda$. Освен това $Po(\lambda)$ се получава като гранично разпределение на $Bi(n, p)$ при $n \to \infty$, $p \to 0$, $np \to \lambda$.

\section{Сравнение на основните разпределения}

\begin{center}
\begin{tabular}{|l|c|c|c|c|}
  \hline
  \textbf{Разпределение} & \textbf{Носител} & \textbf{Параметри} & $\mathbb{E}X$ & $\mathbb{D}X$ \\
  \hline
  Равномерно $U\{1,\dots,n\}$ & $\{1,\dots,n\}$ & $n$ & $\dfrac{n+1}{2}$ & $\dfrac{n^2-1}{12}$ \\[8pt]
  Биномно $Bi(n,p)$ & $\{0,\dots,n\}$ & $n,\,p$ & $np$ & $npq$ \\[4pt]
  Геометрично $Ge(p)$ & $\{0,1,\dots\}$ & $p$ & $\dfrac{q}{p}$ & $\dfrac{q}{p^2}$ \\[8pt]
  Поасоново $Po(\lambda)$ & $\{0,1,\dots\}$ & $\lambda$ & $\lambda$ & $\lambda$ \\[4pt]
  \hline
\end{tabular}
\end{center}

\end{FlushLeft}

\end{document}